\documentclass[spanish,letterpaper,oneside]{article}
\def\documenttitle{Reflexión pedagógica}
\def\documentsubtitle{Temas selectos de matemáticas I}
\def\documentsubject{Ficha de reflexión}
\def\documentauthor{Mtro. Martín Jonathan de la Cruz Muñoz}
\def\coursename{Temas selectos de matemáticas I}
\def\coursecode{PFM02}
\def\universityname{Instituto de Investigación, Innovación y Estudios de Posgrado para la Educación del Estado de Nuevo León}
\def\universityfaculty{Maestría en Aprendizaje y Enseñanza de las Matemáticas}
\def\universitydepartment{Temas selectos de matemáticas I}
\def\universitydepartmentimage{img/iiiepe/marca-iiiepe.png}
\def\universitydepartmentimagecfg{height=1.75cm}
\def\universitylocation{Monterrey, Nuevo León}
\def\studentfullname{Mtro. Martín Jonathan de la Cruz Muñoz}
\def\studentid{00874}
\def\studentemail{martin.delacruz@campus.iiiepe.edu.mx}
\def\programperiod{1er tetramestre, agosto--noviembre de 2026}
\def\programmodality{Modalidad presencial}
\def\authortable{
  \begin{tabular}{ll}
    \textbf{Alumno:} & \studentfullname \\
    Matrícula: & \studentid \\
    Correo institucional: & \studentemail \\
    Unidad académica: & \coursename \\
    Programa: & \programperiod \\
    Modo: & \programmodality \\
  
  \end{tabular}
}
\begin{document}
\authortable
\section*{Reflexión pedagógica}
\noindent\textbf{Instrucciones:} Responde las siguientes preguntas basadas en tu experiencia en la actividad.
\begin{enumerate}
  \item ¿Qué estrategias pedagógicas aplicaste para fomentar la reflexión crítica en tus estudiantes?
  \item ¿Cómo evaluaste el impacto de estas estrategias en el aprendizaje de tus estudiantes?
  \item ¿Qué cambios realizarías para mejorar la actividad en futuras implementaciones?
\end{enumerate}
\section*{Conclusiones}
\noindent\textbf{Espacio para tus conclusiones:}
\end{document}